\documentclass[10pt]{beamer}
\usetheme{metropolis}
\usepackage{graphicx}
\usepackage{listings}
\usepackage{booktabs}
\usepackage{xcolor}
\usepackage{hyperref}

\definecolor{codebg}{HTML}{F5F5F5}
\definecolor{codegreen}{HTML}{2E7D32}
\definecolor{codegray}{HTML}{757575}
\definecolor{codeblue}{HTML}{1565C0}

\lstdefinestyle{rstyle}{
  language=R, backgroundcolor=\color{codebg}, basicstyle=\ttfamily\scriptsize,
  keywordstyle=\color{codeblue}\bfseries, stringstyle=\color{codegreen},
  commentstyle=\color{codegray}\itshape, breaklines=true, frame=single,
  rulecolor=\color{codegray!30}, showstringspaces=false, tabsize=2,
  morekeywords={library,read_csv,filter,select,mutate,arrange,group_by,
    summarise,count,slice_max,drop_na,replace_na,distinct,
    pivot_longer,pivot_wider,left_join,inner_join,anti_join,
    fct_reorder,fct_recode,as.numeric,as.Date,str_to_lower,
    ggplot,aes,geom_col,geom_point,geom_errorbar,coord_flip,
    theme_minimal,labs,n,mean,sd,sqrt}
}
\lstdefinestyle{pythonstyle}{
  language=Python, backgroundcolor=\color{codebg}, basicstyle=\ttfamily\scriptsize,
  keywordstyle=\color{codeblue}\bfseries, stringstyle=\color{codegreen},
  commentstyle=\color{codegray}\itshape, breaklines=true, frame=single,
  rulecolor=\color{codegray!30}, showstringspaces=false, tabsize=4,
  morekeywords={import,from,as,pd,np,plt,read_csv,query,assign,sort_values,
    groupby,agg,merge,melt,pivot_table,dropna,fillna,drop_duplicates,
    to_numeric,to_datetime,value_counts,reset_index,rename}
}

\title{Tidy Data \& Data Wrangling for Viz}
\subtitle{Workshop 2 --- Module 9}
\author{Prof.Asoc.\ Endri Raco}
\institute{Data Visualization for Data Scientists \\ UNYT Tirana}
\date{2025/26}

\begin{document}

\begin{frame}
  \titlepage
\end{frame}

\begin{frame}{Module Roadmap}
  \begin{enumerate}
    \item Tidy data: Wickham's three rules
    \item Wide vs long format: \texttt{pivot\_longer()} / \texttt{melt()}
    \item Six core dplyr verbs with pandas equivalents
    \item The wrangle-to-viz pipeline
    \item Common data problems and their R/Python fixes
    \item Joining data frames for multi-source visualizations
  \end{enumerate}
  \begin{exampleblock}{Learning Outcome}
    You will reshape any dataset into the tidy format that ggplot2
    and seaborn require, and chain a complete wrangle-to-plot pipeline
    in both R and Python.
  \end{exampleblock}
\end{frame}

\section{Tidy Data}

\begin{frame}{Wickham's Three Rules}
  \begin{center}
    \includegraphics[width=0.85\textwidth]{figures_w02m09/tidy_rules}
  \end{center}
  \begin{alertblock}{Pro-Tip}
    ggplot2 requires tidy data: each aesthetic mapping (\texttt{aes(x=, y=,
    color=)}) expects a \emph{column name}. If a variable is spread across
    multiple columns (wide format), you must reshape before plotting.
  \end{alertblock}
\end{frame}

\begin{frame}{Wide vs Long Format}
  \begin{center}
    \includegraphics[width=0.92\textwidth]{figures_w02m09/wide_vs_long}
  \end{center}
\end{frame}

\begin{frame}{pivot\_longer() $\leftrightarrow$ pivot\_wider()}
  \begin{center}
    \includegraphics[width=0.92\textwidth]{figures_w02m09/pivot_diagram}
  \end{center}
\end{frame}

\begin{frame}[fragile]{Reshaping in R vs Python}
  \begin{columns}[T]
    \begin{column}{0.48\textwidth}
      \textbf{R (tidyr)}
\begin{lstlisting}[style=rstyle]
# Wide -> Long
df_long <- df_wide |>
  pivot_longer(
    cols = `2020`:`2022`,
    names_to = "year",
    values_to = "gdp_growth"
  )

# Long -> Wide
df_wide <- df_long |>
  pivot_wider(
    names_from = year,
    values_from = gdp_growth
  )
\end{lstlisting}
    \end{column}
    \begin{column}{0.48\textwidth}
      \textbf{Python (pandas)}
\begin{lstlisting}[style=pythonstyle]
# Wide -> Long
df_long = df_wide.melt(
    id_vars=["Country"],
    var_name="year",
    value_name="gdp_growth")

# Long -> Wide
df_wide = df_long.pivot_table(
    index="Country",
    columns="year",
    values="gdp_growth"
).reset_index()
\end{lstlisting}
    \end{column}
  \end{columns}
  \vspace{4pt}
  \begin{exampleblock}{Data Insight}
    \textbf{Rule}: if you need \texttt{aes(color = year)} but year is
    spread across columns, you must \texttt{pivot\_longer()} first.
    The test: ``Can I name a single column for this aesthetic?''
  \end{exampleblock}
\end{frame}

\section{Six Core Verbs}

\begin{frame}{dplyr Verbs with pandas Equivalents}
  \begin{center}
    \includegraphics[width=0.88\textwidth]{figures_w02m09/dplyr_verbs}
  \end{center}
\end{frame}

\begin{frame}[fragile]{The Six Verbs in R}
\begin{lstlisting}[style=rstyle]
library(tidyverse)

apps <- read_csv("googleplaystore.csv") |>
  # 1. filter: keep rows
  filter(Type %in% c("Free", "Paid"), !is.na(Rating)) |>
  # 2. select: keep columns
  select(Category, Rating, Reviews, Type, `Content Rating`) |>
  # 3. mutate: transform columns
  mutate(Reviews = as.numeric(Reviews),
         log_reviews = log10(Reviews + 1)) |>
  # 4. arrange: sort
  arrange(desc(Rating))

# 5-6. group_by + summarise: aggregate
summary <- apps |>
  group_by(Category) |>
  summarise(
    n = n(),
    mean_rating = mean(Rating),
    se = sd(Rating) / sqrt(n()),
    .groups = "drop"
  ) |>
  slice_max(n, n = 8) |>
  mutate(Category = fct_reorder(Category, mean_rating))
\end{lstlisting}
\end{frame}

\begin{frame}[fragile]{The Six Verbs in Python}
\begin{lstlisting}[style=pythonstyle]
import pandas as pd
import numpy as np

apps = (pd.read_csv("googleplaystore.csv")
    # 1. filter (query)
    .query("Type in ['Free', 'Paid']")
    .dropna(subset=["Rating"])
    # 2. select
    [["Category", "Rating", "Reviews", "Type", "Content Rating"]]
    # 3. mutate (assign)
    .assign(
        Reviews=lambda d: pd.to_numeric(d["Reviews"], errors="coerce"),
        log_reviews=lambda d: np.log10(d["Reviews"] + 1))
    # 4. arrange (sort_values)
    .sort_values("Rating", ascending=False))

# 5-6. group_by + summarise (groupby + agg)
summary = (apps
    .groupby("Category")
    .agg(n=("Rating", "count"),
         mean_rating=("Rating", "mean"),
         se=("Rating", lambda x: x.std() / np.sqrt(len(x))))
    .reset_index()
    .nlargest(8, "n")
    .sort_values("mean_rating"))
\end{lstlisting}
\end{frame}

\section{The Wrangle-to-Viz Pipeline}

\begin{frame}{read $\rightarrow$ clean $\rightarrow$ aggregate $\rightarrow$ plot}
  \begin{center}
    \includegraphics[width=0.95\textwidth]{figures_w02m09/pipeline}
  \end{center}
\end{frame}

\begin{frame}[fragile]{Complete Pipeline: R}
  \begin{columns}[T]
    \begin{column}{0.52\textwidth}
\begin{lstlisting}[style=rstyle]
read_csv("googleplaystore.csv") |>
  # WRANGLE
  filter(Type %in% c("Free","Paid"),
         !is.na(Rating)) |>
  group_by(Category) |>
  summarise(
    n = n(),
    m = mean(Rating),
    se = sd(Rating)/sqrt(n()),
    .groups = "drop") |>
  slice_max(n, n = 8) |>
  mutate(Category =
    fct_reorder(Category, m)) |>

  # PLOT
  ggplot(aes(x = Category,
             y = m)) +
  geom_col(fill = "#1565C0",
           width = 0.5, alpha = 0.6) +
  geom_errorbar(aes(
    ymin = m - se, ymax = m + se),
    width = 0.2) +
  coord_flip() +
  theme_minimal() +
  labs(title = "Mean Rating (+-SE)",
       x = NULL, y = "Mean Rating")
\end{lstlisting}
    \end{column}
    \begin{column}{0.45\textwidth}
      \includegraphics[width=\textwidth]{figures_w02m09/r_wrangle}

      \vspace{4pt}
      {\small The entire pipeline --- from
      CSV to publication chart --- in
      one continuous chain. The pipe
      \texttt{|>} connects wrangling
      verbs; the \texttt{+} connects
      ggplot2 layers.}
    \end{column}
  \end{columns}
\end{frame}

\begin{frame}[fragile]{Complete Pipeline: Python}
  \begin{columns}[T]
    \begin{column}{0.52\textwidth}
\begin{lstlisting}[style=pythonstyle]
summary = (pd.read_csv(
    "googleplaystore.csv")
  # WRANGLE
  .query("Type in ['Free','Paid']")
  .dropna(subset=["Rating"])
  .groupby("Category")
  .agg(n=("Rating","count"),
       m=("Rating","mean"),
       se=("Rating",
           lambda x: x.std()/
           np.sqrt(len(x))))
  .reset_index()
  .nlargest(8, "n")
  .sort_values("m"))

# PLOT
fig, ax = plt.subplots(figsize=(6,5))
ax.barh(summary["Category"],
    summary["m"],
    color="#2E7D32", height=0.5,
    alpha=0.6, xerr=summary["se"],
    capsize=3, ecolor="#333")
for i, (m, s) in enumerate(zip(
    summary["m"], summary["se"])):
    ax.text(m+s+0.01, i, f"{m:.2f}",
        va="center", fontsize=7)
\end{lstlisting}
    \end{column}
    \begin{column}{0.45\textwidth}
      \includegraphics[width=\textwidth]{figures_w02m09/py_wrangle}

      \vspace{4pt}
      {\small pandas method chaining
      mirrors R's pipe. The chain
      goes from \texttt{read\_csv} through
      \texttt{groupby().agg()} to a
      sorted summary, then plots
      with matplotlib.}
    \end{column}
  \end{columns}
\end{frame}

\section{Common Data Problems}

\begin{frame}{Six Problems and Their Fixes}
  \begin{center}
    \includegraphics[width=0.92\textwidth]{figures_w02m09/data_problems}
  \end{center}
\end{frame}

\begin{frame}[fragile]{Cleaning the Google Play Store Data}
\begin{lstlisting}[style=rstyle]
# R: complete cleaning pipeline
apps_clean <- read_csv("googleplaystore.csv") |>
  # 1. Remove duplicates
  distinct(App, .keep_all = TRUE) |>
  # 2. Fix types
  mutate(
    Reviews = as.numeric(Reviews),
    Installs = as.numeric(str_remove_all(Installs, "[+,]")),
    Price = as.numeric(str_remove(Price, "\\$")),
    Last_Updated = mdy(`Last Updated`)
  ) |>
  # 3. Filter valid rows
  filter(Type %in% c("Free", "Paid"), !is.na(Rating)) |>
  # 4. Fix inconsistent categories
  mutate(Category = str_to_title(str_replace_all(Category, "_", " "))) |>
  # 5. Remove outliers
  filter(between(Rating, 1, 5))
\end{lstlisting}
\end{frame}

\section{Joining Data}

\begin{frame}{Four Join Types}
  \begin{center}
    \includegraphics[width=0.92\textwidth]{figures_w02m09/join_types}
  \end{center}
\end{frame}

\begin{frame}[fragile]{Joins in R vs Python}
  \begin{columns}[T]
    \begin{column}{0.48\textwidth}
      \textbf{R (dplyr)}
\begin{lstlisting}[style=rstyle]
# Inner join
df <- apps |>
  inner_join(categories_meta,
    by = "Category")

# Left join (keep all apps)
df <- apps |>
  left_join(country_codes,
    by = c("Country" = "code"))

# Anti join (find unmatched)
missing <- apps |>
  anti_join(categories_meta,
    by = "Category")
\end{lstlisting}
    \end{column}
    \begin{column}{0.48\textwidth}
      \textbf{Python (pandas)}
\begin{lstlisting}[style=pythonstyle]
# Inner join (merge)
df = apps.merge(
    categories_meta,
    on="Category",
    how="inner")

# Left join
df = apps.merge(
    country_codes,
    left_on="Country",
    right_on="code",
    how="left")

# Anti join (no direct method)
missing = apps[~apps["Category"]
    .isin(categories_meta
          ["Category"])]
\end{lstlisting}
    \end{column}
  \end{columns}
\end{frame}

\section{Synthesis}

\begin{frame}{Wrangle-to-Viz Checklist}
  Before plotting, verify:
  \begin{enumerate}
    \item \textbf{Tidy?} Every variable is a column; every observation a row.
    \item \textbf{Types correct?} Numbers are numeric, dates are dates,
          categories are factors.
    \item \textbf{Missing values?} Handled by \texttt{drop\_na()} or
          \texttt{replace\_na()}.
    \item \textbf{Duplicates?} Removed by \texttt{distinct()}.
    \item \textbf{Reshaped?} Wide data pivoted to long if needed for
          \texttt{aes(color = variable)}.
    \item \textbf{Aggregated?} \texttt{group\_by() + summarise()} computed
          before plotting summaries.
    \item \textbf{Sorted?} Categories reordered by value with
          \texttt{fct\_reorder()}.
  \end{enumerate}
\end{frame}

\begin{frame}{Key Takeaways}
  \begin{enumerate}
    \item \textbf{Tidy data} (one variable per column, one observation per
          row) is the required input format for ggplot2 and seaborn.
    \item \textbf{pivot\_longer()} / \texttt{melt()} converts wide to long;
          \textbf{pivot\_wider()} / \texttt{pivot\_table()} does the reverse.
    \item Six core verbs --- \textbf{filter, select, mutate, arrange,
          group\_by, summarise} --- handle 90\% of wrangling tasks.
    \item The \textbf{pipe} (\texttt{|>} in R, method chaining in pandas)
          connects wrangling to plotting in one continuous flow.
    \item \textbf{Always clean before plotting}: fix types, handle NAs,
          remove duplicates, reshape if needed.
    \item \textbf{Joins} combine multiple data sources; use
          \texttt{left\_join()} by default to preserve all primary records.
  \end{enumerate}
\end{frame}

\begin{frame}{Essential Readings}
  \begin{itemize}
    \item Wickham, H. (2014). ``Tidy Data.'' \emph{Journal of Statistical
          Software}, 59(10), 1--23.
    \item Wickham, H. \& Grolemund, G. (2023). \emph{R for Data Science},
          2nd ed., Chapters 3--7 (Data Transformation, Tidy Data, Joins).
    \item McKinney, W. (2022). \emph{Python for Data Analysis}, 3rd ed.,
          Chapters 7--8 (Data Wrangling, Joining).
  \end{itemize}
  \vspace{6pt}
  \textbf{Next Module}: Comparative Lab --- Grammar in Practice (W02-M10)
\end{frame}

\end{document}
